\begin{lemma}[The matching-sector bond inclusion]%
    \label{lem:peps_matching_sector_bond_inclusion}
    \lean{TNLean.PEPS.endpointEmbeddingMatrix,
        TNLean.PEPS.endpointEmbeddingMatrix_isIsometry,
        TNLean.PEPS.blockBondInclusion,
        TNLean.PEPS.blockBondInclusion_isIsometry,
        TNLean.PEPS.blockBondInclusion_mulVec}
    \leanok
    Let $V=\bigoplus_{i\in I}V_i$ be a finite orthogonal direct sum.
    Any injective relabelling of orthonormal endpoint coordinates defines
    an isometric inclusion. In particular, the natural inclusion
    \begin{align}
        E_V:\bigoplus_{i\in I}(V_i\otimes V_i)
            \longrightarrow V\otimes V
    \end{align}
    is an isometry. If a bond vector is obtained by arranging the entries
    of matrices $X_i$ in their respective summands, its image under $E_V$
    is the vector of entries of $\bigoplus_i X_i$.
\end{lemma}

\begin{proof}\leanok
    Each basis vector maps to a distinct pair of endpoint basis vectors.
    Entries in different summands vanish; entries in a common summand
    are unchanged.
\end{proof}

\begin{definition}[Multiplicity restoration in endpoint coordinates]%
    \label{def:peps_multiplicity_block_bond_matrix}
    \lean{TNLean.PEPS.multiplicityBlockBondMatrix,
        TNLean.PEPS.multiplicityBlockBondMatrix_isIsometry,
        TNLean.PEPS.multiplicityBlockBondMatrix_mulVec_sqrt}
    \leanok
    \uses{def:peps_multiplicity_bond_coefficients}
    Choose nonempty finite multiplicity sets $M_i$, with $m_i=|M_i|$.
    In paired endpoint coordinates the bond isometry is
    \begin{align}
        J\ket{i,a,b}=\frac{1}{\sqrt{m_i}}
            \sum_{u\in M_i}\ket{i,(a,u),(b,u)}.
    \end{align}
    Thus $J$ carries the vector of entries of $\sqrt{m_i}X_i$ to the
    vector of entries of $X_i\otimes I_{M_i}$, block by block.
    This is the bond transformation of
    \cite[Section~7]{Schuch2010PEPS}.
\end{definition}

\begin{proof}\leanok
    Relabel the coefficient isometry by the two restored endpoints.
    Its square-root normalization cancels the bond weight.
\end{proof}

\begin{theorem}[The bond map on the full endpoint spaces]%
    \label{thm:peps_full_multiplicity_bond_map}
    \lean{TNLean.PEPS.fullMultiplicityBondMap,
        TNLean.PEPS.fullMultiplicityBondMap_initialProjection,
        TNLean.PEPS.fullMultiplicityBondMap_mulVec_sqrt_blockDiagonal,
        TNLean.PEPS.fullMultiplicityBondMap_dotProduct}
    \leanok
    \uses{lem:peps_matching_sector_bond_inclusion,
        def:peps_multiplicity_block_bond_matrix}
    Put $V'=\bigoplus_i(V_i\otimes\mathbb C^{M_i})$ and define
    \begin{align}
        F=E_{V'}JE_V^\dagger:V\otimes V\longrightarrow V'\otimes V'.
    \end{align}
    Then $F^\dagger F=E_VE_V^\dagger$, so $F$ preserves inner products
    on the matching-sector subspace. More precisely, whenever
    $E_VE_V^\dagger\phi=\phi$,
    \begin{align}
        \langle F\psi,F\phi\rangle=\langle\psi,\phi\rangle.
    \end{align}
    It sends the full endpoint vector of
    $\bigoplus_i\sqrt{m_i}X_i$ to that of
    $\bigoplus_i(X_i\otimes I_{M_i})$.
    For $m_i=\dim V_i$, this is the local physical bond transformation
    used in \cite[Section~7]{Schuch2010PEPS}.
\end{theorem}

\begin{proof}\leanok
    Both inclusions and $J$ are isometries. Cancelling their Gram
    operators gives $F^\dagger F=E_VE_V^\dagger$. The overlap identity
    follows by moving $F$ to the other side of the inner product.
    On a block-diagonal bond vector, $E_V^\dagger$ recovers its
    matching-sector coordinates, to which the preceding action of $J$
    applies.
\end{proof}
